\bookchapter{The finite handoff to computation}
\label{sec:vii-handoff}

\subsection{What leaves the matter volume}

The matter volume ends with objects that a computation can receive without changing their scientific status. They include typed particle rows, native generator forms, admission and stability fields, mass-lift channels, nuclear population ledgers, element and isotope addresses, carrier budgets, closure addresses and selected relation receipts. Some are exact algebraic constructions; some are catalog records; some carry measured comparators or fitted parameters. A faithful computation preserves those distinctions.

The receiving volume develops the Substrate Ledger Computer and its language, exact-state representations, write lifecycle and applications. It does not need to repeat every physical argument in order to serialize a matter row. It does need to know what the row means, which operations are defined on it and which results remain only conditional or empirical. The handoff is therefore a contract between types and operators, rather than a claim that execution supplies missing physical evidence.\cite{viisama,viispine}

The conceptual route is finite matter grammar, typed mass and binding channels, carrier/container return, computation grammar, language and write lifecycle. This route connects the volumes while leaving ownership of each result clear. Matter definitions and comparisons remain in Volume II. Runtime behavior and computational exploitation belong to Volume III. A successful execution may validate its implementation of the received rule; it does not by itself validate the rule as a description of nature.

\subsection{The row as a serializable scientific object}

The row identity introduced in Chapter~\ref{sec:vii-foundations} contains a source table, schema, address and acting operator. A computational representation must retain enough of that identity to distinguish rows whose numerical fields happen to agree. It must also retain the difference between an input and a derived field.

For example, a native value, a debit and a resulting readout are three related fields. If the representation stores only the readout, a null result could be indistinguishable from a row that never had a native candidate. The LC06 witness shows why the history matters: native value $18$ and debit $18$ produce a null readout, and that cancellation has a different meaning from an absent record. Exact arithmetic can preserve the cancellation only if the operands and their types remain available.

Likewise, a known particle or element label should remain a reveal field. Moving the label into the address-generation stage would alter the inference order even if the displayed table remained unchanged. The computational interface must preserve which information was available at construction, which was used for selection and which was opened only for comparison.

The same applies to admission. A nonzero native value is not a replacement for the admission predicate. A rejected candidate can be represented exactly and can remain valuable as a control. A hidden-support row can be needed for the grammar while being barred from matter-source emission. The receiving language should be able to carry those statuses without turning every stored record into an admitted physical object.

This is not a new implementation specification. It states the semantic content of the source objects already used in the volume. Different implementations may encode the fields differently, but they cannot omit a load-bearing distinction and still claim to reproduce the same scientific object.

\subsection{Exact rationals before decimal display}

Several central quantities are rational:
\begin{equation}
\frac18,\qquad\frac5{48},\qquad\frac{17}{192},
\end{equation}
\begin{equation}
\kappa=\frac{7117}{768},\qquad
\frac{7093}{192}.
\end{equation}
The exact forms support cancellation and comparison without introducing a display tolerance into the underlying identity. A decimal representation is useful for communicating scale, but it should not replace the rational source when the claim concerns exact equality.

The normalized-square asymmetry relation is an example. Its coefficient follows from the exact difference between two source channels. Storing a rounded decimal for each channel and subtracting them can introduce a residual that belongs to the representation rather than the model. Retaining the rational forms makes the source of exactness explicit and separates it from the precision used for later measured comparisons.

Exact representation does not remove the need for units. The rational $7117/768$ is a structural coefficient in the declared channel. A measurement in MeV or atomic mass units enters through its own conversion and comparison specification. Exactness of the rational coefficient does not turn that conversion into a derived physical law.

Undefined expressions require equal care. The $K=X=0$ carrier row remains part of the lift table while $Y=X/K$ is undefined. A receiving program that fills the quotient with zero would preserve the table's length but change its mathematical content. Filling it with one would impose a different unsupported convention. The appropriate representation distinguishes a retained row with an inapplicable or undefined quotient from an ordinary numerical value.

This distinction also affects aggregate statistics. A reported closure of 138 evaluable quotients is not a failure to account for a 139-row table. It is a statement that one row lies outside the quotient's domain. A computation that silently includes an invented endpoint could appear more complete while actually weakening the specification.

\subsection{Source inventories and executable grammars}

The executable N100 basis contains one hundred typed computation rows. It is not a replacement census for the physical-source surfaces. The matter volume has already used 35 compact mass rows, 139 lift-table rows, 321 catalog rows, 126 inventory rows, an 80-row stable structural surface and a 126-row periodic surface. Their counts remain attached to their own selection predicates.

A computation may prune or select a basis for an operation. The cross-volume source gives the example of removing five $p{:}9$ rows from an executable basis. Such a selection does not delete rows from the particle catalog, the mass table or the periodic source. It changes the operation's admitted domain. The source inventories must remain recoverable without applying the computation's pruning as a scientific revision.

This is analogous to the distinction between a table and a query over the table, but the consequences here are mathematical. An executable basis can omit a source row because the chosen operator is not defined on it, because it is redundant for a particular representation, or because the operation uses a narrower class. None of those possibilities is equivalent to declaring that the source row is physically nonexistent or experimentally rejected.

A program can therefore consume a 126-row periodic input while using a 100-row operation grammar. It can compare a 139-row lift table while evaluating only 138 quotients. It can preserve a 321-row catalog that includes carriers and rejected candidates while selecting a smaller matter-allowed subset for source support. Correctness requires the selection rule and its domain, not equality of all counts.

The source-authorized N100-to-F81 elementwise join remains absent in the inspected record. Its aggregate contact is not enough to fill the missing assignment. A receiving computation can carry both source and target objects and can record that the join is undefined; it need not invent the map to remain a complete representation of the current state.

\subsection{Writes receipts and retained history}

The W8/X1/W9 lifecycle supplies a computation-facing language for a supported relation, activation and completion. In the current typed reading, these roles are distinct from homogeneous mass-lift tensors. Theta18 supplies a carrier or completion-report role at the relevant interface, while the A-field carrier supplies transport. The matter volume exports these meanings rather than merely the displayed values $8$, $1$, $9$ and $18$.

The reception results show why an endpoint value can be an insufficient record. Different relation histories or incomplete delivery patterns can preserve an aggregate number while changing the reconstructed object. A computation that retains only the final count would lose the information needed to distinguish those cases. Exact receipts preserve relation identity, direction, address and payload at the scope established by the source test.

The AFC2 and AFC3 chain is a concrete example. A scalar coordinate such as $u=d=9$ does not identify the six-object incidence structure. The eighteen target-barred records carry information that the aggregate omits. The geometry is reconstructed from that information under a declared procedure. The subsequent delivery test shows that a retained symmetry can survive some losses even though the exact shape changes. A successful serialization must therefore preserve more than the diagnostic that happens to remain constant.

The relation between an informational receipt and a physical event remains a separate interpretive question. An exact reversible record can demonstrate fidelity within a finite computational model. It does not automatically establish a physical energy, a measured length or a causal matter-generation process. Keeping the receipt exact is nevertheless valuable: it preserves the construction so that its actual content can be inspected without relying on a picture or a compressed summary.

\subsection{Empirical fields retain their provenance}

The matter account includes several kinds of empirical contact. A neutrino absolute scale uses a measured splitting. A historical nuclear readout uses fitted coefficients and then evaluates a frozen extension. B5T identifies an informational tensor before binding targets are opened, then evaluates a binding residual lane with its own development and held data. A target-visible He-4 discovery calculation uses information that a blind confirmation would withhold.

These histories remain part of the exported object. A numerical prediction cannot be interpreted solely from its final residual. The receiving program or report must preserve which quantities were fixed, which were fit, which targets were visible, which rows were held and which controls were part of the declared test. The same residual can mean different things under different information histories.

A useful example is the phrase ``zero refit.'' For CR277 it identifies the extension stage: the earlier fitted operator is frozen and applied to a broader representative table. It does not say that the original coefficients were never fit. Similarly, a typed-rational coefficient selected near a fitted value retains that selection history even when the final displayed coefficient is exact. Exact arithmetic and empirical selection are compatible, but their combination is not the same as a coefficient derived without data contact.

The B5T target-blind tensor stage also must be stated precisely. Freezing the assignment before binding is opened prevents binding from selecting that assignment or breaking its tie. It does not mean that every subsequent component of the binding prediction is unsupervised or that the measured scale is derived from a finite permutation alone. The representation and the residual test have related but distinct evidence burdens.

An executable record is most useful when it preserves these distinctions rather than replacing them with a single PASS flag. The flag identifies the outcome of declared predicates. The predicates, inputs, controls and classification explain what that outcome means.

\subsection{Failures are part of the finite export}

Several unsuccessful propositions are load-bearing parts of this volume. The direct source-to-rest-mass identity failed its locked control. CR241 retained low holdout residual while failing candidate uniqueness. The proposed linear constituent binding account did not supply the nonlinear binding readout. The two CR283 geometry searches returned no valid maps. Frozen hydrogen-ladder energy transfers missed their measured targets. Target-visible refinements remain discoveries rather than confirmations.

A computational export that retains only successful rows would make the source record less informative. The failures distinguish quantities that would otherwise be easy to merge. The direct mass failure separates coupling from rest-source information. The failed geometry searches separate cardinality from placement. The energy-transfer misses separate exact structural graphs from a calibrated physical law.

Failures also need their scope. A failed specific map has the scope of the tested mapping. A failed frozen energy expression does not invalidate every relation in the graph on which it was evaluated. Conversely, a surviving graph or exact identity does not rescue the rejected physical prediction. The export must retain the proposition tested, not merely a favorable or unfavorable label.

This principle avoids two opposite distortions. One would erase useful exact work whenever a stronger interpretation fails. The other would use useful exact work to imply that the stronger interpretation succeeded. The finite record supports neither shortcut. It permits each result to retain the status earned by its own evidence.

\subsection{Source custody for the completed account}

The continuous matter synthesis in SAMA supplies the historical narrative foundations. The expanded chapters receive the later completed source reports, derivations and current authority summarized in the source register. Its frozen repository version preserves the source chapters and the copied Courtroom material. The SAM Research Project technical spine supplies the source-bound recovery map, while its live binding and ATOM3D documents control the current status of those branches. The evidence-rich unapproved review assembly supplies focused derivation and deviation chains without conferring author approval by its existence.\cite{viisama,viispine,viireview,viibinding,viiatom}

The sources have different authority roles. A later prose assembly does not automatically supersede an executed result. A repository reorganization does not change a scientific classification. A historical result remains evidence for its own test even when a later domain authority supplies the current synthesis. Where a source identifies a missing permanent test key, the missing key remains a custody boundary rather than being replaced by a similar identifier.

The numerical receipts quoted in this book are reported from those records unless the text explicitly describes code inspection. The preparation of this manuscript did not rerun scientific campaigns, refit coefficients or reproduce every independent validator. Selected code inspection establishes what the inspected routine does, such as retaining a fitted coefficient when a typed substitution falls outside its threshold. It does not substitute for execution of every dependency or independent replication of the entire research program.

This distinction does not reduce a source-based book to a mere list of claims. The exact algebra can be followed directly, the definitions can be stated without ambiguity, and the historical tests can be described with their controls and limitations. The resulting account is complete at the level it declares: an organized mathematical construction and its recorded evidence, with physical identifications bounded by the evidence actually available.

\subsection{The matter account at its stopping point}

The dedicated spin section adds exact joint state counting, controlled thermal ordering, microscopic memory, composition cancellation, source-specific kinetic clocks and the latest paid physical response to the matter construction. These pass to Volume III as source-bound exact factors, generators, certified matrix operands and retained dependency closures.

The work developed here has a definite finite content. It supplies a typed distinction between substrate support and admitted matter; complementary carrier and retained channels; several explicitly different particle and matter surfaces; class-dependent generator and lift laws; conditional Higgs and neutrino packets; a four-slot nuclear ledger; an exact channel-gap asymmetry relation; historical fitted and later held binding readouts; periodic-family and isotope organization; finite reception and reconstruction results; complete conditional nuclear and thermal sources; causal ordered contact histories; occupied Pauli atoms; full-motion molecular recipients; isotope recoil transfer; a held light-source test; and an explicit common physical energy convention.

It also has definite limits. Candidate uniqueness does not follow from a small residual. The absolute neutrino scale retains its external anchor. A periodic frontier identity is not an observation. The completed atomic and molecular recipients retain their declared source and prepared-family scopes. An informational tensor or unscaled graph embedding is not yet a calibrated physical nuclear law. The later chapters identify the achieved source, field, occupied-state and outgoing-record maps, together with the specific physical assignments still open. These statements are the boundary of the present account, not a schedule of obligations appended to it.

Volume III receives the objects on the established side of that boundary together with the statuses that protect their meaning. It can preserve them, transform them under declared operators and examine their computational consequences. It need not solve every open physical question before doing so. The finite handoff is thus the appropriate end of the matter volume: exact structure where the construction is exact, recorded comparison where the evidence is empirical, and explicit limits where the source has not completed the bridge.
